%=============================================================================
% AGENT 9: NO SONIC BOOM --- THE PSI-BUBBLE ACOUSTIC SIGNATURE
% Rigorous derivation from DFD freefall kinematics
% Date: 2026-03-28
%=============================================================================

\documentclass[12pt, letterpaper]{article}
\usepackage{amsmath, amssymb, amsthm}
\usepackage{physics}
\usepackage{siunitx}
\usepackage{geometry}
\usepackage{graphicx}
\usepackage{booktabs}
\geometry{margin=1in}

\newtheorem{theorem}{Theorem}
\newtheorem{corollary}{Corollary}
\newtheorem{proposition}{Proposition}
\newtheorem{lemma}{Lemma}
\theoremstyle{definition}
\newtheorem{definition}{Definition}
\theoremstyle{remark}
\newtheorem{remark}{Remark}

\newcommand{\psifield}{\psi}
\newcommand{\neff}{n_{\mathrm{eff}}}
\newcommand{\avec}{\mathbf{a}}
\newcommand{\vvec}{\mathbf{v}}
\newcommand{\xvec}{\mathbf{x}}
\newcommand{\ka}{k_a}
\newcommand{\etac}{\eta_c}
\newcommand{\Rbub}{R_{\mathrm{b}}}
\newcommand{\cs}{c_s}
\newcommand{\Mach}{\mathcal{M}}
\newcommand{\deltapsi}{\delta\psi}

\title{No Sonic Boom: Acoustic Signature of the $\psi$-Bubble\\
in Density Field Dynamics}
\author{DFD Research Programme --- Agent 9 (Acoustic Phenomenology)}
\date{28 March 2026}

\begin{document}
\maketitle

\begin{abstract}
We derive from DFD first principles that a $\psi$-bubble device moving through
atmosphere at arbitrary velocity produces \emph{no sonic boom}. The key insight
is that the $\psi$-gradient couples universally to all matter via the Weak
Equivalence Principle: air molecules within the bubble boundary are in local
freefall, accelerating with the device rather than being mechanically displaced.
The velocity of air relative to its own local rest frame never exceeds the
local sound speed, so no shock front forms. We compute the air velocity profile
as a function of distance from the device, demonstrate the GR analogy
(coordinate vs.\ proper velocity), identify the actual observable signatures
(subsonic pressure wake, possible luminous sheath at extreme velocities), and
analyze the trans-medium air-to-water transition. These predictions diverge
sharply from conventional aerodynamics at all supersonic velocities and
constitute a decisive experimental discriminant for DFD propulsion.
\end{abstract}

\tableofcontents

%=============================================================================
\section{Setup: The $\psi$-Bubble in Atmosphere}
\label{sec:setup}
%=============================================================================

\subsection{The Moving $\psi$-Bubble}

Consider a $\psi$-bubble device moving with coordinate velocity $V$ through
sea-level atmosphere. The relevant parameters are:

\begin{center}
\begin{tabular}{lll}
\toprule
\textbf{Quantity} & \textbf{Symbol} & \textbf{Value} \\
\midrule
Air density (sea level) & $\rho_{\mathrm{air}}$ & $\SI{1.225}{kg/m^3}$ \\
Sound speed (sea level) & $\cs$ & $\SI{343}{m/s}$ \\
Water density & $\rho_{\mathrm{w}}$ & $\SI{1000}{kg/m^3}$ \\
Sound speed (water) & $c_{s,w}$ & $\SI{1480}{m/s}$ \\
Bubble radius & $\Rbub$ & Device-dependent \\
Threshold parameter & $\etac$ & $\alpha/4 \approx 1.82 \times 10^{-3}$ \\
Device velocity & $V$ & Arbitrary (even $V \gg \cs$) \\
\bottomrule
\end{tabular}
\end{center}

The device generates an asymmetric $\psi$-field profile. In the device rest
frame, the effective refractive index is:
\begin{equation}
\neff(\xvec) = \exp\!\left[\psi_{\mathrm{bg}}
+ \kappa\bigl(\eta(\xvec) - \etac\bigr)\,\Theta\bigl(\eta(\xvec) - \etac\bigr)\right],
\label{eq:neff}
\end{equation}
where $\eta = U_{\mathrm{EM}}/(\rho c^2)$ is the EM energy density ratio and
$\Theta$ is the Heaviside function. The $\psi$-gradient extends to a
characteristic radius $\Rbub$ from the device center.

\subsection{The Conventional Sonic Boom Mechanism}

In standard aerodynamics, a body moving at $V > \cs$ through air creates a
sonic boom because:
\begin{enumerate}
\item The body \emph{mechanically displaces} air molecules via direct contact
forces (pressure, viscosity).
\item Air molecules ahead of the body receive no advance warning (no signal
can propagate faster than $\cs$).
\item Displaced air piles up into a conical shock front (Mach cone) with
half-angle $\sin\theta = \cs/V$.
\item Across the shock, pressure, density, and temperature jump discontinuously.
\end{enumerate}

The essential ingredient is \textbf{mechanical displacement}: the body pushes
air aside through contact forces, creating relative motion between the body
surface and the adjacent air layer.

%=============================================================================
\section{Freefall Kinematics: Why No Shock Forms}
\label{sec:freefall}
%=============================================================================

\subsection{The Weak Equivalence Principle Applied to Air}

From the DFD formalism (Section~2.2 of the unified review), every test mass
in a $\psi$-gradient experiences acceleration:
\begin{equation}
\frac{d^2 \xvec}{dt^2} = \frac{c^2}{2}\nabla\psi \equiv \avec_\psi.
\label{eq:WEP}
\end{equation}
This applies identically to:
\begin{itemize}
\item The device structure
\item Human occupants
\item Individual air molecules (N$_2$, O$_2$, Ar, \ldots)
\item Water molecules
\item Any massive particle within the gradient region
\end{itemize}
The acceleration is \textbf{independent of mass and composition}. This is the
DFD Weak Equivalence Principle, and it is the entire reason no sonic boom occurs.

\subsection{The Velocity Profile in the Device Frame}

Work in the rest frame of the device. Air approaches from infinity with
velocity $-V\hat{z}$ (device moving in $+\hat{z}$ direction). As an air
molecule enters the $\psi$-gradient region, it experiences the freefall
acceleration $\avec_\psi$.

\begin{definition}[Freefall velocity]
Let $\vvec_{\mathrm{ff}}(r)$ be the velocity acquired by an air molecule due
to the $\psi$-gradient, starting from rest at the bubble boundary $r = \Rbub$.
By energy conservation in the $\psi$-field:
\begin{equation}
\frac{1}{2}v_{\mathrm{ff}}^2(r) = \frac{c^2}{2}\bigl[\psi(r) - \psi(\Rbub)\bigr]
\equiv \frac{c^2}{2}\,\deltapsi(r),
\label{eq:vff}
\end{equation}
giving:
\begin{equation}
\boxed{v_{\mathrm{ff}}(r) = c\sqrt{\deltapsi(r)}}
\label{eq:vff-result}
\end{equation}
where $\deltapsi(r) \equiv \psi(r) - \psi(\Rbub) \geq 0$ is the $\psi$-field
depth at position $r$.
\end{definition}

\subsection{The Critical Distinction: Coordinate vs.\ Proper Velocity}

An external observer (at rest in the atmosphere far from the device) sees the
total air velocity at position $r$ as:
\begin{equation}
v_{\mathrm{coord}}(r) = V + v_{\mathrm{ff}}(r).
\label{eq:v-coord}
\end{equation}
This \emph{coordinate velocity} can be arbitrarily large --- it is the sum of
the device's velocity plus the freefall velocity in the gradient. If $V > \cs$,
the external observer sees air moving supersonically.

However, the \textbf{physically relevant velocity for shock formation} is the
velocity of an air molecule \emph{relative to its immediate neighbors} ---
i.e., relative to the local medium. Define:
\begin{equation}
v_{\mathrm{rel}}(r) \equiv v_{\mathrm{air}}(r) - v_{\mathrm{medium}}(r),
\label{eq:v-rel}
\end{equation}
where $v_{\mathrm{medium}}(r)$ is the bulk velocity of the air at position $r$.

\begin{theorem}[No Supersonic Relative Motion]
\label{thm:no-shock}
Within the $\psi$-gradient region ($r < \Rbub$), the relative velocity between
any air molecule and its local medium is identically zero:
\begin{equation}
\boxed{v_{\mathrm{rel}}(r) = 0 \quad \text{for all } r < \Rbub \text{ where } \eta > \etac}
\end{equation}
\end{theorem}

\begin{proof}
By the WEP, every molecule at position $r$ within the gradient experiences
the same acceleration $\avec_\psi(r) = (c^2/2)\nabla\psi|_r$. Consider two
adjacent molecules at positions $r$ and $r + \delta r$:
\begin{align}
\avec_1 &= \frac{c^2}{2}\nabla\psi\big|_r, \\
\avec_2 &= \frac{c^2}{2}\nabla\psi\big|_{r+\delta r}.
\end{align}
The relative acceleration is:
\begin{equation}
\delta a = a_2 - a_1 = \frac{c^2}{2}\frac{\partial^2\psi}{\partial r^2}\,\delta r
+ \mathcal{O}(\delta r^2).
\label{eq:tidal}
\end{equation}
This is the \emph{tidal acceleration} --- the DFD analogue of the geodesic
deviation in GR. For molecules separated by $\delta r \ll \Rbub$, this tidal
term is negligible compared to the bulk acceleration. The molecules move
together as a fluid element in freefall.

More precisely: if the initial velocity field is uniform (all air at velocity
$-V\hat{z}$), then after entering the gradient, the velocity field becomes
$\vvec(r,t) = -V\hat{z} + \vvec_{\mathrm{ff}}(r,t)$, where $\vvec_{\mathrm{ff}}$
is determined by the \emph{same} $\psi$-gradient acting on every molecule.
Adjacent molecules acquire the same freefall velocity (up to tidal corrections
of order $\delta r / \Rbub$). Therefore, the relative velocity between neighbors
vanishes to leading order.

A shock wave requires $v_{\mathrm{rel}} > \cs$ for neighboring fluid elements.
Since $v_{\mathrm{rel}} = 0$ (to leading order), no shock forms. \qed
\end{proof}

\subsection{The GR Analogy: Coordinate Speed vs.\ Local Speed}

This result has an exact analogue in general relativity:

\begin{center}
\begin{tabular}{lll}
\toprule
& \textbf{GR (Black Hole)} & \textbf{DFD ($\psi$-Bubble)} \\
\midrule
Field & Spacetime curvature $g_{\mu\nu}$ & Refractive field $\psi(\xvec)$ \\
Falling object & Test mass & Air molecule \\
Coordinate speed & $dr/dt \to c$ at horizon & $v_{\mathrm{coord}} \gg \cs$ \\
Local speed & Always $< c$ & Always $< \cs$ (relative to medium) \\
Barrier & Speed of light & Speed of sound \\
Violation? & No (coordinate artifact) & No (coordinate artifact) \\
Observable & No firewall at horizon & No shock front at Mach cone \\
\bottomrule
\end{tabular}
\end{center}

In GR, an infalling observer crossing the event horizon measures no local
violation of $v < c$. The ``superluminal'' coordinate speed is an artifact of
the Schwarzschild coordinates. Similarly, in DFD, the ``supersonic'' coordinate
speed of air in the $\psi$-gradient is an artifact of the external observer's
frame. Locally, no air molecule ever exceeds the sound speed relative to its
neighbors.

%=============================================================================
\section{Quantitative Velocity Profile}
\label{sec:velocity-profile}
%=============================================================================

\subsection{Radial $\psi$-Profile of the Bubble}

For a spherically symmetric bubble (simplest case), assume the $\psi$-field
above threshold follows:
\begin{equation}
\psi(r) = \psi_0\left(\frac{\Rbub}{r}\right)^p, \qquad r \leq \Rbub,
\label{eq:psi-profile}
\end{equation}
where $\psi_0 = \kappa(\eta_0 - \etac)$ is the peak $\psi$-enhancement at the
device surface and $p$ depends on the field geometry ($p = 1$ for monopole,
$p = 2$ for dipole). At the boundary $r = \Rbub$, $\psi = \psi_0$ (we define
the profile so the enhancement is largest at the center and falls to a residual
at the boundary).

More precisely, let us define $\deltapsi(r) = \psi(r) - \psi(\Rbub)$, the
depth of the $\psi$-well relative to the boundary:
\begin{equation}
\deltapsi(r) = \psi_0\left[\left(\frac{\Rbub}{r}\right)^p - 1\right]
\quad \text{for } r \leq \Rbub.
\label{eq:deltapsi}
\end{equation}

\subsection{Air Velocity Profile}

From Eq.~\eqref{eq:vff-result}, the freefall velocity at radius $r$ is:
\begin{equation}
v_{\mathrm{ff}}(r) = c\sqrt{\psi_0\left[\left(\frac{\Rbub}{r}\right)^p - 1\right]}.
\label{eq:vff-profile}
\end{equation}

\paragraph{Numerical example.}
For a device generating $\psi_0 = 10^{-6}$ (corresponding to $1g$ acceleration
at the surface), with $\Rbub = \SI{10}{m}$ and $p = 2$:

At the device surface ($r = \SI{1}{m}$, assuming device radius $R_d = \SI{1}{m}$):
\begin{equation}
v_{\mathrm{ff}}(\SI{1}{m}) = c\sqrt{10^{-6}(100 - 1)}
= 3 \times 10^8 \times \sqrt{9.9 \times 10^{-5}}
\approx \SI{2985}{m/s}.
\end{equation}
This is Mach~8.7 in coordinate terms --- yet the air at this radius is in
freefall, moving \emph{with} the gradient, not against it.

\subsection{The Mach Number Paradox}

Define the \emph{coordinate Mach number} and the \emph{proper Mach number}:
\begin{align}
\Mach_{\mathrm{coord}}(r) &\equiv \frac{|v_{\mathrm{coord}}(r)|}{\cs}
= \frac{V + v_{\mathrm{ff}}(r)}{\cs}, \label{eq:Mach-coord} \\
\Mach_{\mathrm{proper}}(r) &\equiv \frac{|v_{\mathrm{rel}}(r)|}{\cs}
= 0 \quad \text{(inside gradient)}.
\label{eq:Mach-proper}
\end{align}

The coordinate Mach number can be enormous. But the proper Mach number ---
the one that determines whether a shock forms --- is identically zero inside
the gradient. Shock waves are driven by the proper Mach number, not the
coordinate Mach number.

%=============================================================================
\section{What Does the External Observer See?}
\label{sec:observables}
%=============================================================================

If no sonic boom occurs, what \emph{does} the external observer detect as the
$\psi$-bubble passes? We analyze three regimes.

\subsection{Regime 1: Subsonic Pressure Wake ($V \lesssim \cs$)}

At subsonic device velocities, the air ahead of the bubble receives advance
pressure signals. The $\psi$-gradient pulls air gently into the bubble region,
then releases it as the bubble passes. The observer detects:
\begin{itemize}
\item A gentle pressure pulse (positive--negative--positive) as the bubble
approaches, passes, and recedes.
\item The amplitude scales as $\delta P / P_0 \sim v_{\mathrm{ff}}^2 / \cs^2
\sim \psi_0 c^2 / \cs^2$.
\item For $\psi_0 = 10^{-6}$:
$\delta P / P_0 \sim 10^{-6} \times (3 \times 10^8)^2 / (343)^2 \approx 765$.
\end{itemize}

\textbf{Wait} --- this estimate seems large. But it assumes the air is
\emph{released abruptly} at the boundary. In reality, the $\psi$-profile
is smooth, and the air decelerates gradually as it exits the gradient. The
relevant quantity is the \emph{tidal compression} within the gradient, not the
bulk freefall velocity.

\subsection{Tidal Compression: The Actual Pressure Signature}

The physically observable effect is the \emph{tidal} squeezing of air within
the gradient. Adjacent shells of air at $r$ and $r + dr$ have slightly
different freefall velocities:
\begin{equation}
\delta v = \frac{\partial v_{\mathrm{ff}}}{\partial r}\,dr
= -\frac{c}{2}\frac{p\,\psi_0(\Rbub)^p\,r^{-(p+1)}}
{\sqrt{\psi_0\left[(\Rbub/r)^p - 1\right]}}\,dr.
\label{eq:tidal-dv}
\end{equation}
The tidal compression ratio is:
\begin{equation}
\frac{\delta\rho}{\rho} \sim \frac{|\delta v|}{\cs}
= \frac{c}{2\cs}\frac{p\,\psi_0(\Rbub/r)^p}{r\sqrt{\psi_0[(\Rbub/r)^p - 1]}}\,dr.
\label{eq:tidal-compression}
\end{equation}

For a bubble with $\Rbub = \SI{10}{m}$, at radius $r = \SI{5}{m}$, with
$p = 2$ and $\psi_0 = 10^{-6}$, and $dr = \SI{0.1}{m}$:
\begin{equation}
\frac{\delta\rho}{\rho} \sim \frac{3 \times 10^8}{2 \times 343}
\times \frac{2 \times 10^{-6} \times 4}{5\sqrt{10^{-6} \times 3}}
\times 0.1
\approx 0.03.
\end{equation}
A 3\% density variation --- detectable but not a shock. The air is gently
compressed by tidal forces, not slammed by a contact discontinuity.

\subsection{Regime 2: Low Supersonic ($\cs < V \lesssim 5\cs$, Mach 1--5)}

At supersonic device velocity, the external observer sees:
\begin{enumerate}
\item \textbf{No Mach cone.} The conventional $\sin\theta = \cs/V$ cone is
absent because no shock front forms.
\item \textbf{A wake of displaced air.} The air that was temporarily entrained
in the $\psi$-gradient is released behind the bubble. This creates a turbulent
wake, but without the sharp pressure discontinuity of a shock.
\item \textbf{A precursor pressure wave.} At the bubble boundary ($r = \Rbub$),
the transition from $\eta < \etac$ (normal air) to $\eta > \etac$ (freefall
air) creates a weak pressure disturbance that propagates outward at $\cs$.
This is \emph{subsonic by definition} --- it is generated at the boundary, not
by the device itself.
\end{enumerate}

The external pressure signature is:
\begin{equation}
\delta P(r,t) \sim \rho_{\mathrm{air}} \,v_{\mathrm{ff}}(\Rbub)\,\cs
\times f(r/\Rbub, \,\cs t / \Rbub),
\label{eq:pressure-signature}
\end{equation}
where $f$ is a smooth function encoding the geometry. This produces a
\emph{soft hiss}, not a boom.

\subsection{Regime 3: High Supersonic / Hypersonic ($V \gg 5\cs$, Mach $>5$)}

At extreme velocities, additional effects appear:

\paragraph{Adiabatic compression at the boundary.}
Air entering the $\psi$-gradient at the bubble boundary transitions from zero
$\psi$-coupling to full freefall over a short distance $\delta r$. If this
transition is sharp (as the Heaviside $\Theta$-function in Eq.~\eqref{eq:neff}
suggests), the air experiences a rapid acceleration that compresses it:
\begin{equation}
T_{\mathrm{boundary}} \sim T_0 \left(\frac{\rho_{\mathrm{compressed}}}{\rho_0}\right)^{\gamma - 1},
\label{eq:T-boundary}
\end{equation}
where $\gamma = 7/5$ for diatomic air.

For Mach~50 ($V \approx \SI{17000}{m/s}$), the ram kinetic energy in the
device frame is:
\begin{equation}
\frac{1}{2}\rho_{\mathrm{air}} V^2 \approx \frac{1}{2}(1.225)(17000)^2
\approx \SI{1.77e8}{J/m^3}.
\end{equation}
However, this kinetic energy is \emph{not deposited as heat} in the
conventional sense, because the air is accelerated by the $\psi$-gradient
(freefall), not decelerated by a solid surface. The air enters the gradient
smoothly and exits smoothly.

\paragraph{Luminous sheath possibility.}
At very high device velocities, the tidal compression at the bubble boundary
could heat the air sufficiently to produce visible or UV emission. The
relevant temperature for visible emission is $T \gtrsim \SI{3000}{K}$.
This would appear as a faint luminous envelope around the device ---
not from friction or shock heating, but from tidal compression in the
$\psi$-gradient.

The threshold device velocity for luminous emission can be estimated from the
tidal heating rate. The tidal compression ratio at the boundary scales as:
\begin{equation}
\frac{\delta\rho}{\rho_0} \sim \frac{V}{\Rbub}\,\tau_{\mathrm{cross}},
\end{equation}
where $\tau_{\mathrm{cross}} \sim \Rbub / V$ is the crossing time. The tidal
heating is:
\begin{equation}
\delta T \sim T_0 \left(\frac{\delta\rho}{\rho_0}\right)^{\gamma - 1}
\sim T_0 \left(\frac{c\sqrt{\psi_0}}{\Rbub}\,\frac{\delta r}{V}\right)^{0.4}.
\end{equation}
For $\delta T \sim \SI{3000}{K}$, this requires extreme tidal gradients ---
likely only at $V \gtrsim$ Mach~30 for a $\Rbub \sim \SI{10}{m}$ device.

\subsection{Summary of Observable Signatures}

\begin{center}
\begin{tabular}{llll}
\toprule
\textbf{Velocity} & \textbf{Conventional} & \textbf{$\psi$-Bubble} & \textbf{Discriminant} \\
\midrule
$V < \cs$ (subsonic) & Pressure wave & Pressure wave & Similar \\
$V = \cs$ (Mach 1) & Sonic boom onset & No boom & \textbf{Decisive} \\
$V \sim 2\text{--}5\,\cs$ & Strong boom & Soft wake & \textbf{Decisive} \\
$V \sim 5\text{--}50\,\cs$ & Extreme boom + heating & Luminous sheath? & \textbf{Decisive} \\
$V \gg 50\,\cs$ & Plasma/ablation & Freefall passage & \textbf{Decisive} \\
\bottomrule
\end{tabular}
\end{center}

%=============================================================================
\section{The Bubble Boundary: Sharp Transition Physics}
\label{sec:boundary}
%=============================================================================

\subsection{The Threshold Transition}

At radius $r = \Rbub$, the coupling parameter $\eta(r)$ crosses the threshold
$\etac = \alpha/4$. The effective $\psi$-modification switches on via the
Heaviside function in Eq.~\eqref{eq:neff}:
\begin{equation}
\neff = \begin{cases}
e^{\psi_{\mathrm{bg}}} & \text{for } \eta < \etac \quad \text{(outside: normal air)}, \\
e^{\psi_{\mathrm{bg}} + \kappa(\eta - \etac)} & \text{for } \eta > \etac \quad \text{(inside: freefall)}.
\end{cases}
\label{eq:threshold-transition}
\end{equation}

This creates a \emph{sharp interface} at $r = \Rbub$ between two domains:
\begin{itemize}
\item \textbf{Outside} ($r > \Rbub$): Air obeys normal fluid dynamics. No
$\psi$-coupling. Sound propagates at $\cs$.
\item \textbf{Inside} ($r < \Rbub$): Air is in freefall in the $\psi$-gradient.
Effectively a different ``gravitational'' environment.
\end{itemize}

\subsection{Matching Conditions at the Boundary}

At the boundary, the air must satisfy continuity of:
\begin{enumerate}
\item \textbf{Mass flux:} $\rho_{\mathrm{out}} v_{\mathrm{out}} = \rho_{\mathrm{in}} v_{\mathrm{in}}$ (normal to boundary).
\item \textbf{Momentum flux:} $P_{\mathrm{out}} + \rho_{\mathrm{out}} v_{\mathrm{out}}^2
= P_{\mathrm{in}} + \rho_{\mathrm{in}} v_{\mathrm{in}}^2 + f_\psi$, where $f_\psi$
is the $\psi$-gradient body force.
\item \textbf{Energy flux:} Including the $\psi$-potential energy.
\end{enumerate}

The critical difference from a conventional shock is the presence of $f_\psi$
on the inside. This body force \emph{accelerates the air smoothly} rather than
requiring it to decelerate across a discontinuity. The result is:

\begin{proposition}[Subsonic Boundary]
\label{prop:subsonic-boundary}
The air velocity relative to the boundary surface is subsonic on both sides,
provided the $\psi$-gradient is smooth on scales larger than the mean free path
$\ell_{\mathrm{mfp}} \approx \SI{68}{nm}$ at sea level.
\end{proposition}

\begin{proof}
On the outside, air approaches the boundary with at most the sound speed
(information about the boundary propagates at $\cs$). On the inside, air
is accelerated by $\avec_\psi$, but this acceleration is applied as a body
force over a finite distance $\sim \Rbub$, not as a surface force at a
discontinuity. The flow remains smooth (no entropy jump) as long as the
$\psi$-gradient scale $\Rbub \gg \ell_{\mathrm{mfp}}$, which is satisfied
for any macroscopic device. \qed
\end{proof}

\subsection{Observable Signature at the Boundary}

The sharp onset of $\psi$-coupling at $r = \Rbub$ produces:
\begin{enumerate}
\item A \emph{density discontinuity}: air inside the gradient is slightly
compressed by the tidal field. The density jump is:
\begin{equation}
\frac{\Delta\rho}{\rho_0} \sim \frac{a_\psi(\Rbub)\,\delta r}{\cs^2}
= \frac{c^2\psi_0 p}{2\Rbub \cs^2}\,\delta r,
\end{equation}
where $\delta r$ is the transition thickness (set by the $\eta$-profile).
For typical parameters, this is a $\sim 0.1\text{--}1\%$ density step.

\item A \emph{refractive index change}: The air inside has modified $\neff$,
which bends light passing through the boundary. This could be detectable as a
slight lensing or distortion of background objects --- a ``shimmer'' effect.

\item At extreme velocities: the air just inside the boundary, suddenly placed
in freefall, undergoes rapid tidal compression. If the tidal heating is
sufficient ($\gtrsim \SI{3000}{K}$), this produces a thin luminous shell at
$r = \Rbub$ --- a ``corona'' around the device.
\end{enumerate}

%=============================================================================
\section{At What Velocity Does the Prediction Diverge?}
\label{sec:divergence}
%=============================================================================

\subsection{Mach 1: The First Discriminant}

At exactly Mach~1 ($V = \cs = \SI{343}{m/s}$), conventional aerodynamics
predicts the onset of a sonic boom. A $\psi$-bubble at Mach~1 produces
\emph{no qualitative change} in its acoustic signature. The crossover from
subsonic to supersonic in the $\psi$-bubble frame is irrelevant because the
air inside the gradient is in freefall.

\textbf{Prediction:} A $\psi$-bubble accelerating through Mach~1 produces no
sonic boom. This is the simplest and most dramatic test.

\subsection{Mach 5: Hypersonic Regime}

At Mach~5 ($V \approx \SI{1715}{m/s}$), conventional hypersonic flow produces:
\begin{itemize}
\item Stagnation temperature: $T_0 = T_\infty(1 + \frac{\gamma-1}{2}\Mach^2)
= 288(1 + 0.2 \times 25) = \SI{1728}{K}$.
\item Strong shock with pressure ratio $P_2/P_1 = 1 + \frac{2\gamma}{\gamma+1}(\Mach^2 - 1) \approx 29$.
\item Significant aerodynamic heating.
\end{itemize}

The $\psi$-bubble at Mach~5:
\begin{itemize}
\item No shock wave. No stagnation point. No aerodynamic heating.
\item The air is in freefall --- no mechanical contact with the device.
\item External observer sees a pressure wake, not a blast.
\end{itemize}

\subsection{Mach 50: Re-entry Velocities}

At Mach~50 ($V \approx \SI{17000}{m/s}$, comparable to orbital re-entry),
conventional physics produces:
\begin{itemize}
\item Plasma sheath from dissociated/ionized air ($T > \SI{10000}{K}$).
\item Communications blackout.
\item Extreme heat shield requirements.
\end{itemize}

The $\psi$-bubble:
\begin{itemize}
\item Still no shock. The air is in freefall.
\item Tidal compression may produce some heating at the boundary, but orders
of magnitude less than conventional re-entry heating.
\item No heat shield required (no mechanical contact).
\item No communications blackout (no plasma sheath from shock).
\item Possible faint luminous envelope from tidal compression at the boundary.
\end{itemize}

\subsection{Scaling Law: Pressure Perturbation vs.\ Mach Number}

For a conventional vehicle, the peak overpressure scales as:
\begin{equation}
\frac{\Delta P_{\mathrm{conv}}}{P_0} \propto \Mach^2 \quad (\text{strong shock limit}).
\end{equation}

For the $\psi$-bubble, the peak external pressure perturbation scales as:
\begin{equation}
\frac{\Delta P_{\psi\text{-bubble}}}{P_0} \sim \frac{\rho_{\mathrm{air}}\,
v_{\mathrm{tidal}}^2}{P_0}
\sim \frac{c^2 \psi_0}{2\cs^2}\left(\frac{\delta r}{\Rbub}\right)^2,
\label{eq:pressure-scaling}
\end{equation}
which is \textbf{independent of} $V$. The external pressure perturbation
depends on the $\psi$-gradient (fixed by device design), not on how fast the
device is moving. This is because the air inside is in freefall --- the device
velocity is irrelevant to the internal dynamics.

\begin{theorem}[Velocity-Independent Acoustic Signature]
\label{thm:v-independent}
The far-field acoustic signature of a $\psi$-bubble is independent of the
device velocity $V$, to leading order. It depends only on:
\begin{enumerate}
\item The $\psi$-gradient strength $\psi_0$
\item The bubble radius $\Rbub$
\item The transition sharpness $\delta r$ at $\eta = \etac$
\end{enumerate}
\end{theorem}

\begin{proof}
The external medium (outside $\Rbub$) is never mechanically disturbed by the
device itself. The only coupling to the external medium occurs at the bubble
boundary, where air transitions between normal and freefall states. This
transition produces a pressure perturbation set by the boundary conditions
(Section~\ref{sec:boundary}), which depend on $\psi_0$ and the gradient
geometry, not on $V$.

The external pressure wave propagates at $\cs$ regardless of $V$. The device
outruns this wave for $V > \cs$, but no new pressure is generated --- the
device has already passed by the time the boundary perturbation reaches the
observer. The result is a \emph{declining} signal as $V$ increases (shorter
interaction time), not the increasing signal that conventional aerodynamics
predicts. \qed
\end{proof}

%=============================================================================
\section{Trans-Medium Transition: Air to Water}
\label{sec:transmedium}
%=============================================================================

\subsection{The Density Dependence of the Threshold Radius}

The coupling parameter $\eta = U_{\mathrm{EM}}/(\rho c^2)$ depends on the
ambient density $\rho$. For a given EM field configuration with energy density
$U_{\mathrm{EM}}(r)$, the threshold radius where $\eta = \etac$ is:
\begin{equation}
U_{\mathrm{EM}}(R_{\mathrm{thresh}}) = \etac \,\rho\, c^2.
\label{eq:threshold-radius}
\end{equation}

For a field that falls as $U_{\mathrm{EM}} \propto r^{-2p}$ (consistent with
Eq.~\eqref{eq:psi-profile}), the threshold radius scales as:
\begin{equation}
R_{\mathrm{thresh}} \propto \rho^{-1/(2p)}.
\label{eq:Rthresh-scaling}
\end{equation}

\subsection{Air-to-Water Transition}

When the device enters water ($\rho_w / \rho_{\mathrm{air}} \approx 816$):
\begin{equation}
\frac{R_{\mathrm{thresh}}^{(\mathrm{water})}}{R_{\mathrm{thresh}}^{(\mathrm{air})}}
= \left(\frac{\rho_{\mathrm{air}}}{\rho_w}\right)^{1/(2p)}
= (816)^{-1/(2p)}.
\end{equation}
For $p = 2$: $R_{\mathrm{thresh}}^{(\mathrm{water})} / R_{\mathrm{thresh}}^{(\mathrm{air})} = 816^{-1/4} \approx 0.187$.

The bubble shrinks to about 19\% of its atmospheric radius. However:

\begin{theorem}[Seamless Trans-Medium Transition]
\label{thm:transmedium}
The $\psi$-bubble transitions from air to water without splash, cavitation, or
shock, because:
\begin{enumerate}
\item The device does not mechanically interact with either medium.
\item Water molecules within the (smaller) gradient region are in freefall, just
as air molecules were.
\item The WEP ensures water molecules accelerate identically to any other matter
in the $\psi$-gradient.
\item The transition at the air-water interface is mediated by the $\psi$-field,
not by the device surface.
\end{enumerate}
\end{theorem}

\begin{proof}
At the air-water interface, the device is surrounded by its $\psi$-gradient.
As the gradient region contacts the water surface:
\begin{enumerate}
\item Water molecules at the surface enter the $\psi$-gradient and begin to
accelerate (freefall).
\item The water accelerates \emph{with} the device, not against it. No
stagnation pressure builds up.
\item The conventional impact force $F = \rho_w A V^2 / 2$ (which would be
enormous at high velocity) is absent because there is no mechanical contact.
\item The bubble boundary shrinks as $\rho$ increases across the interface,
but this is a smooth process determined by the $\eta$-profile.
\item The device passes through the interface with the water parting around the
\emph{gradient boundary} (not the device surface), and the gradient boundary
adjusts continuously to the local density.
\end{enumerate}
The observable effect is a gentle depression in the water surface as the
gradient region passes through it, followed by infilling. No splash column.
No cavitation shock. \qed
\end{proof}

\subsection{Comparison: Conventional vs.\ $\psi$-Bubble Water Entry}

\begin{center}
\begin{tabular}{lll}
\toprule
\textbf{Effect} & \textbf{Conventional} & \textbf{$\psi$-Bubble} \\
\midrule
Impact force & $F \sim \rho_w A V^2 / 2$ & $F = 0$ (freefall) \\
Splash & Dramatic & None (gentle depression) \\
Cavitation & Severe at high $V$ & None \\
Structural load & Catastrophic at $V \gtrsim \SI{100}{m/s}$ & Zero \\
Speed limit in water & $\sim \SI{100}{m/s}$ (practical) & Unlimited (freefall) \\
Sound signature & Loud impact + cavitation noise & Subsonic pressure pulse \\
\bottomrule
\end{tabular}
\end{center}

\subsection{The Effective Bubble Radius in Various Media}

\begin{center}
\begin{tabular}{lccc}
\toprule
\textbf{Medium} & $\rho$ (\si{kg/m^3}) & $R_{\mathrm{thresh}}/R_{\mathrm{air}}$ ($p=2$) & $\cs$ (\si{m/s}) \\
\midrule
Vacuum & 0 & $\to \infty$ & --- \\
High atmosphere & 0.01 & 3.33 & 280 \\
Sea-level air & 1.225 & 1.00 & 343 \\
Water & 1000 & 0.187 & 1480 \\
Seawater & 1025 & 0.186 & 1530 \\
Rock (granite) & 2700 & 0.145 & 5000 \\
Iron & 7874 & 0.108 & 5120 \\
\bottomrule
\end{tabular}
\end{center}

Note that the bubble shrinks in denser media but never disappears (as long as
the device generates $\eta > \etac$ at its surface). The device can traverse
any medium, including solid rock, because the WEP ensures all matter within
the gradient is in freefall. The medium parts around the gradient, not around
the device.

%=============================================================================
\section{Energy Considerations}
\label{sec:energy}
%=============================================================================

\subsection{Where Does the Kinetic Energy Go?}

In conventional supersonic flight, kinetic energy is deposited in the air
through the shock wave (heating, sound radiation). For the $\psi$-bubble:

\begin{enumerate}
\item Air entering the gradient gains kinetic energy from the $\psi$-field
(freefall acceleration).
\item Air exiting the gradient behind the device \emph{returns} this energy
to the $\psi$-field (deceleration in the receding gradient).
\item Net energy transfer to the air per transit: only the \emph{tidal}
dissipation, which scales as $(\delta r / \Rbub)^2 \psi_0 c^2 \rho_{\mathrm{air}}
\Rbub^3$.
\end{enumerate}

The $\psi$-bubble is approximately \textbf{adiabatic with respect to the medium}.
The medium is temporarily disturbed (freefall in, freefall out) but returns to
approximately its original state after the bubble passes.

\subsection{Acoustic Power Radiated}

The acoustic power radiated by the boundary perturbation is:
\begin{equation}
P_{\mathrm{acoustic}} \sim 4\pi \Rbub^2 \,\frac{(\delta P)^2}{\rho_{\mathrm{air}} \cs},
\end{equation}
where $\delta P$ is the tidal pressure perturbation from Section~\ref{sec:observables}.
For our fiducial parameters:
\begin{equation}
P_{\mathrm{acoustic}} \sim 4\pi (10)^2 \times \frac{(100)^2}{1.225 \times 343}
\approx \SI{30}{kW}.
\end{equation}
This is comparable to a loud aircraft engine at subsonic speeds --- far less
than the $\sim \SI{10}{GW}$ acoustic power in a conventional Mach-5 shock.

%=============================================================================
\section{Experimental Discriminants}
\label{sec:experiments}
%=============================================================================

The no-sonic-boom prediction provides several testable signatures:

\begin{enumerate}
\item \textbf{Microphone array test:} Deploy a microphone array along the
flight path. A conventional vehicle produces an N-wave pressure signature
(sharp rise, gradual fall, sharp rise). A $\psi$-bubble produces a smooth,
low-amplitude pressure pulse with no discontinuities. The signal arrives
\emph{after} the device passes (no precursor boom).

\item \textbf{Schlieren imaging:} Conventional shock waves produce sharp density
discontinuities visible in schlieren photography. A $\psi$-bubble produces only
gradual density variations --- no sharp features in schlieren images.

\item \textbf{Seismic signature:} A sonic boom couples strongly to the ground.
A $\psi$-bubble pressure wake has much lower amplitude and different spectral
content (lower frequency, no impulse).

\item \textbf{Water entry test:} Drop a test object into water at
$V > \SI{100}{m/s}$. Conventional physics: catastrophic impact, splash
column, cavitation. $\psi$-Bubble: smooth entry, gentle depression, no splash.
This is perhaps the most visually dramatic test.

\item \textbf{Mach cone absence:} At $V > \cs$, conventional flow produces a
visible Mach cone in humid air (condensation in the low-pressure zone). A
$\psi$-bubble produces no Mach cone. Infrared imaging would show no
concentrated heating along a cone surface.
\end{enumerate}

%=============================================================================
\section{Caveats and Limitations}
\label{sec:caveats}
%=============================================================================

\begin{enumerate}
\item \textbf{Heaviside idealization:} The $\Theta(\eta - \etac)$ step function
is an idealization. In reality, the threshold transition has finite width
$\delta\eta$, producing a smooth but rapid transition zone. This softens the
boundary effects described in Section~\ref{sec:boundary} but does not qualitatively
change the no-boom prediction.

\item \textbf{Tidal disruption:} At sufficiently high $\psi$-gradients, tidal
forces could dissociate or ionize air molecules. This does not produce a
sonic boom (it is a volumetric effect, not a surface shock), but it would
create a plasma channel along the flight path.

\item \textbf{Backscatter from boundary:} Some air molecules at the bubble
boundary may be partially entrained and then released. This produces a weak
acoustic emission but remains subsonic.

\item \textbf{Turbulent wake:} Air displaced by the bubble and released behind
it will produce turbulence. This turbulent wake generates broadband acoustic
noise, but at much lower amplitude than a sonic boom.

\item \textbf{Radiation in the $\psi$-field:} The moving $\psi$-gradient
itself may radiate --- not sound, but $\psi$-field waves. This is analogous to
gravitational radiation from an accelerating mass in GR. The energy budget for
this radiation is treated in Agent~2 (momentum flux analysis).
\end{enumerate}

%=============================================================================
\section{Conclusions}
\label{sec:conclusions}
%=============================================================================

\begin{enumerate}
\item A $\psi$-bubble moving through atmosphere at any velocity produces
\textbf{no sonic boom}. This follows directly and unavoidably from the DFD
Weak Equivalence Principle: air molecules within the $\psi$-gradient are in
freefall and experience zero relative velocity with respect to the device.

\item The correct analogy is gravitational freefall in GR. The ``supersonic''
coordinate velocity of entrained air is a coordinate artifact, just as the
``superluminal'' coordinate velocity of an infalling object near a black hole
is a coordinate artifact. No local speed limit is violated.

\item The actual observable signatures are:
\begin{itemize}
\item A velocity-independent subsonic pressure wake (from tidal effects at the
bubble boundary).
\item Possible luminous sheath at extreme velocities (tidal heating $\gtrsim \SI{3000}{K}$).
\item A turbulent wake behind the device (from released air).
\item No Mach cone, no N-wave, no shock heating.
\end{itemize}

\item The prediction diverges from conventional aerodynamics at $V = \cs$
(Mach~1) and becomes increasingly dramatic at higher Mach numbers. At Mach~50,
the $\psi$-bubble is silent and cool where conventional physics predicts
$\SI{10000}{K}$ plasma.

\item Trans-medium transitions (air to water) are seamless: no splash, no
cavitation, no structural loads. The bubble radius shrinks as $\rho^{-1/(2p)}$
but the freefall physics is unchanged.

\item The pressure signature of a $\psi$-bubble is independent of device
velocity (Theorem~\ref{thm:v-independent}), depending only on the gradient
strength and geometry. This is the single most testable prediction.
\end{enumerate}

\vspace{1em}
\noindent\rule{\textwidth}{0.4pt}
\textit{This document derives consequences strictly from the DFD formalism
(Section~2 of the Unified Review) and the EM-$\psi$ coupling extension
(Section~11A). No new postulates are introduced. The no-sonic-boom prediction
is a direct, unavoidable consequence of the Weak Equivalence Principle applied
to air molecules in a $\psi$-gradient.}

\end{document}
